\documentclass[a4paper,12pt]{article}
\usepackage[utf8]{inputenc}
\usepackage[T1]{fontenc}
\usepackage[ngerman]{babel}
\usepackage{amsmath}
\usepackage{amssymb}
\usepackage{enumitem}
\usepackage{geometry}
\geometry{a4paper, margin=2.5cm}

\title{Einsendeaufgaben -- Lektion 4\\[0.5em]
\large Modul 61111: Mathematische Grundlagen}
\author{}
\date{}

\begin{document}
\maketitle

\section*{Aufgabe 4.5}

\begin{enumerate}[label=\alph*)]

    \item
    \underline{„$\Rightarrow$"-Richtung:}

    Sei $(a_n)$ und $(b_n)$ konvergent, dann folgt aus Proposition 13.4.12 direkt, dass
    $(a_n + b_n)$ und $(a_n - b_n)$ konvergieren.

    \underline{„$\Leftarrow$"-Richtung:}

    Sei $(a_n + b_n)$ und $(a_n - b_n)$ konvergent.

    Es gilt:
    \[
        a_n = \frac{1}{2}(2a_n) = \frac{1}{2}\bigl((a_n + b_n) + (a_n - b_n)\bigr)
    \]
    \[
        b_n = \frac{1}{2}(2b_n) = \frac{1}{2}\bigl((a_n + b_n) - (a_n - b_n)\bigr)
    \]

    Da $(a_n + b_n)$ und $(a_n - b_n)$ konvergieren, folgt aus Proposition 13.4.12,
    dass $(a_n)$ und $(b_n)$ ebenfalls konvergieren.

    \item
    Gilt nicht, da z.\,B. $(a_n)$ mit $a_n = n$ und $(b_n)$ mit $b_n = -n$ divergieren,
    aber $(a_n + b_n)$ konvergiert, da $a_n + b_n = 0$ für alle $n \in \mathbb{N}$.

    \item
    Gilt nicht.

    Beispiel:
    \[
        a_n = \frac{1}{n}, \quad b_n = (-1)^n
    \]

    Da $(a_n)$ eine Nullfolge und $(b_n)$ beschränkt ist, gilt wegen der Proposition 13.4.9,
    dass $(a_n b_n)$ konvergiert. Allerdings divergiert $(b_n)$.

\end{enumerate}

\end{document}
